%==============================================================================
% FOREWORD
%==============================================================================

\chapter*{Foreword}
\addcontentsline{toc}{chapter}{Foreword}

This is a monograph on what biological systems do with charge.

The argument is simple. Wherever there is charge in a bounded system, it must
be redistributed. Biology is what charge redistribution looks like when it is
sustained, structured, and self-maintaining over geological timescales. The
genome, the membrane, the metabolic network, the immune system, the disease
state --- these are all configurations through which charge is moved, balanced,
and re-balanced. The information storage attributed to nucleic acids, the
catalytic specificity attributed to enzymes, the recognition attributed to
immunoglobulins --- these are downstream consequences of the prior requirement
that charge be handled.

This priority is not a metaphor. The existence cost of any biological
structure --- the thermodynamic price of partitioning matter into a distinct
entity --- is paid in units of charge separation. Entropy, in this framework,
is not a measure of disorder but of the charge state of a partition boundary.
The arrow of time is not mysterious: it points in the direction that partition
boundaries can propagate without reversing the work already done. Maxwell's
Demon fails not because measurement costs information, but because the
topology of charge-locked networks is independent of particle velocity. These
are not analogies; they are derivations.

The framework presented here follows from a single axiom --- that physical
systems occupy bounded regions of phase space admitting partition and nesting
--- and applies it without modification to biological systems at every scale.
The mathematics is self-contained: nothing is assumed beyond what is reviewed
in Appendix~A, and every result is derived from the preceding ones, with no
free parameters and no auxiliary hypotheses. The reader who finds a numerical
prediction will find its derivation two pages earlier and its empirical
validation two pages after.

The book is structured in seven parts. Part~I develops the formal language.
Part~II uses it to dissolve four problems that have resisted solution for
decades. Parts~III through~VI apply it progressively to catalysis, the genome,
oscillatory networks, and disease. Part~VII demonstrates that nucleotide
sequence analysis is a branch of wave physics and was always so. The synthesis
chapter (Chapter~28) traces the full arc from the two opening axioms to all
principal results.

The parts are written to be accessible in isolation. A clinician working on the
equations of state for disease can enter at Part~VI without first absorbing the
categorical mechanics. A structural biologist interested in nucleic acid charge
dynamics can read Part~IV independently. The full structure rewards reading in
sequence, but the individual results stand on their own.

One terminological note. The word \textit{partition} is used throughout with a
precise technical meaning: a division of phase space into non-overlapping
regions separated by a boundary. It is not used loosely as a synonym for
division, separation, or component. The reader who holds this definition fixed
will find that everything else follows.

\bigskip
\noindent
\textit{Munich, \the\year}

\vspace{0.3cm}
\noindent
K.\,F.\,S.
